\documentclass[11pt]{article}
\usepackage[margin=1in]{geometry}
\usepackage[T1]{fontenc}
\usepackage{lmodern}
\usepackage{amsmath}
\usepackage[numbers,square]{natbib}
\usepackage{cleveref}

\title{Numeric Natbib with Manual Bibliography}
\author{Test Corpus}
\date{}

\begin{document}
\maketitle

\begin{abstract}
This fixture exercises natbib in numeric mode with a hand-written
\texttt{thebibliography} list rather than a BibTeX run, so it compiles
self-contained. It mixes parenthetical and textual citations, a citation with a
page locator, and a multiple-key citation, all of which must resolve to bracketed
numbers in document order.
\end{abstract}

\section{Citation Forms}
\label{sec:forms}

The event-study estimator used throughout applied corporate finance traces to
\citet{brown1985using}, who formalize the market-model abnormal return. Standard
treatments of the cross-sectional test statistics appear in later work
\citep{campbell1997econometrics}. When trading is thin, a trade-to-trade
adjustment is preferable \citep[p.~5]{maynesrumsey1993}, and the small-sample
behaviour of the associated tests has been examined directly
\citep{kolari2010event,corrado1992nonparametric}.

Because the bibliography is configured with the \texttt{numbers} and
\texttt{square} options, each of these citations renders as a bracketed number,
and \citet{brown1985using} additionally prints the author name before its number.
The multiple-key citation above should collapse to a single bracket listing two
numbers.

\section{Cross-References}
\label{sec:xref}

\Cref{sec:forms} introduces the four citation forms, while the present section
confirms that \texttt{cleveref} and \texttt{natbib} coexist. A trivial
identity, $(a+b)^2 = a^2 + 2ab + b^2$, is included only to verify that inline
mathematics is unaffected by the citation machinery. The reference list that
follows is enumerated manually, so the citation numbers are assigned by the order
of the \verb|\bibitem| entries below.

\begin{thebibliography}{9}

\bibitem[Brown and Warner(1985)]{brown1985using}
Stephen J. Brown and Jerold B. Warner.
\newblock Using daily stock returns: The case of event studies.
\newblock \emph{Journal of Financial Economics}, 14(1):3--31, 1985.

\bibitem[Campbell et al.(1997)]{campbell1997econometrics}
John Y. Campbell, Andrew W. Lo, and A. Craig MacKinlay.
\newblock \emph{The Econometrics of Financial Markets}.
\newblock Princeton University Press, Princeton, NJ, 1997.

\bibitem[Maynes and Rumsey(1993)]{maynesrumsey1993}
Elizabeth Maynes and John Rumsey.
\newblock Conducting event studies with thinly traded stocks.
\newblock \emph{Journal of Banking \& Finance}, 17(1):145--157, 1993.

\bibitem[Corrado and Zivney(1992)]{corrado1992nonparametric}
Charles J. Corrado and Terry L. Zivney.
\newblock The specification and power of the sign test in event study hypothesis
tests using daily stock returns.
\newblock \emph{Journal of Financial and Quantitative Analysis}, 27(3):465--478,
1992.

\bibitem[Kolari and Pynn\"{o}nen(2010)]{kolari2010event}
James W. Kolari and Seppo Pynn\"{o}nen.
\newblock Event study testing with cross-sectional correlation of abnormal
returns.
\newblock \emph{Review of Financial Studies}, 23(11):3996--4025, 2010.

\end{thebibliography}

\end{document}
